\section{Flow Matching 与扩散模型的统一视角}

\subsection{等价性的建立}

在本讲的最后部分，讲者揭示了 Flow Matching 和扩散模型之间的深层联系。当我们选择与 VP/VE 调度对应的 $\mu_t$ 和 $\sigma_t$ 时：

\begin{enumerate}
    \item Flow Matching 的条件概率路径 $p_t(x | x_1) = \mathcal{N}(\mu_t x_1, \sigma_t^2 I)$ 与扩散模型的前向跳跃分布完全一致（只是时间方向相反）
    \item Flow Matching 的条件向量场与 Probability Flow ODE 中的漂移项存在精确对应关系
    \item 两种训练目标在适当的参数化和损失权重下是等价的
\end{enumerate}

\begin{importantbox}{统一视角：两条路通向同一座山}
\textbf{扩散模型的路径}：定义前向过程（噪声注入） $\to$ 推导反向 SDE $\to$ 得到 Probability Flow ODE $\to$ 训练 score function

\textbf{Flow Matching 的路径}：定义条件流映射 $\to$ 推导条件向量场 $\to$ 训练向量场预测网络

两条路最终到达同一个目标：一个能将噪声映射到数据的 ODE。不同之处在于：
\begin{itemize}
    \item 扩散模型从概率密度路径出发推导向量场
    \item Flow Matching 从流映射出发推导概率密度路径
    \item 推导顺序相反，但结果等价
\end{itemize}
\end{importantbox}

\subsection{Flow Matching 的优势}

尽管数学上等价，Flow Matching 框架具有以下实践优势：

\begin{table}[H]
\centering
\begin{tabular}{lcc}
\toprule
\textbf{特性} & \textbf{扩散模型} & \textbf{Flow Matching} \\
\midrule
参考分布 & 必须是标准高斯 & 可以是任意分布 \\
训练目标 & 噪声预测 / score 预测 & 速度预测 \\
理论推导 & 需要 SDE/Fokker-Planck & 直接定义 ODE \\
噪声调度 & 需要精心设计 & 自由选择 $\mu_t$, $\sigma_t$ \\
直线路径 & 非自然（弯曲轨迹） & 自然（OT 路径） \\
\bottomrule
\end{tabular}
\caption{扩散模型与 Flow Matching 的特性对比}
\end{table}

\begin{knowledgebox}{为什么现代生成模型更倾向于 Flow Matching}
Stable Diffusion 3（Stability AI）和 Flux（Black Forest Labs）等最新图像生成模型都采用了 Flow Matching / Rectified Flow 框架。主要原因包括：
\begin{enumerate}
    \item \textbf{更直的轨迹}：OT 路径产生的轨迹更接近直线，使得少步采样的效果更好
    \item \textbf{更简单的训练}：速度预测目标比噪声预测或 score 预测更直观
    \item \textbf{更灵活的设计}：可以自由选择参考分布和路径参数化
    \item \textbf{更好的理论基础}：与 optimal transport 理论的天然联系
\end{enumerate}
\end{knowledgebox}

\subsection{本章小结}

Flow Matching 和扩散模型在数学上是等价的，但 Flow Matching 提供了更简洁的理论框架和更灵活的模型设计空间。在实践中，直线路径（OT 路径）带来的更优轨迹结构使得 Flow Matching 成为当前主流生成模型的首选训练范式。


